Two features that the reform left unchanged matter as much as its changes. First, the obligation to
commission a DPE before marketing a property predates it. The requirement, and the liability
attached to marketing a property without the required diagnostic, were established by \emph{loi
ALUR} (Article L271-4 of the \emph{Code de la construction et de l'habitation}) well before 2021.
The legal exposure a seller incurs by coming to market without a certificate is therefore governed
by a mandate that straddles the reform date unchanged, while what July 2021 altered is how much a
\emph{disclosed} grade is worth to the buyer. Second, the reform sharpened a pre-existing
regulatory threshold. Under \emph{loi Climat et R\'esilience}, dwellings
rated E, F, or G face a phased prohibition on new lettings. This draws a bright line between the
bottom of the grade distribution and the rest, and anchors the binary quality distinction used below
in a concrete regulatory threshold.
